% !TEX program = xelatex
\documentclass[10pt,aspectratio=169]{beamer}
\usetheme{Madrid}
\usecolortheme{dolphin}
\usepackage{ctex}
\setCJKmainfont{Noto Sans CJK SC}
\setmainfont{Liberation Sans}

\usepackage{amsmath,amssymb,bm,mathtools}
\usepackage{booktabs,tabularx,array,multirow}
\usepackage{tikz}
\usetikzlibrary{arrows.meta,positioning,calc,fit,shapes.geometric}
\usepackage{graphicx,hyperref}
\usefonttheme{professionalfonts}
\setbeamertemplate{navigation symbols}{}
\setbeamertemplate{footline}[frame number]
\setbeamertemplate{equation number}[default]
\setbeamerfont{frametitle}{size=\large}
\setbeamersize{text margin left=6mm,text margin right=6mm}
\setlength{\parskip}{2pt}\linespread{1.02}
\definecolor{deepblue}{RGB}{27,67,121}
\definecolor{softblue}{RGB}{236,243,252}
\definecolor{softgreen}{RGB}{235,247,241}
\definecolor{softorange}{RGB}{252,243,233}
\definecolor{softred}{RGB}{253,239,238}
\definecolor{deepgreen}{RGB}{32,105,80}
\definecolor{deeporange}{RGB}{174,87,20}
\definecolor{deepred}{RGB}{174,50,45}
\definecolor{purple}{RGB}{93,74,151}
\setbeamercolor{block title}{fg=white,bg=deepblue}
\setbeamercolor{block body}{bg=softblue,fg=black}
\setbeamercolor{block title example}{fg=white,bg=deepgreen}
\setbeamercolor{block body example}{fg=black,bg=softgreen}
\setbeamercolor{block title alerted}{fg=white,bg=deeporange}
\setbeamercolor{block body alerted}{fg=black,bg=softorange}
\setbeamercolor{alerted text}{fg=deeporange}
\newcommand{\R}{\mathbb R}
\newcommand{\E}{\mathbb E}
\newcommand{\dd}{\,\mathrm d}
\newcommand{\Cov}{\operatorname{Cov}}
\newcommand{\N}{\mathcal N}
\newcommand{\Lt}{L_t}
\newcommand{\Mtheta}{M_k^\theta}
\newcommand{\Mref}{M_k^\mathrm{ref}}
\newcommand{\hstar}{\Delta\tau_k^{\star}}
\newcommand{\Rstar}{R_k^{\star}}
\newcommand{\wtilde}{\widetilde w}
\newcommand{\KL}{\operatorname{KL}}
\newcommand{\ESS}{\operatorname{ESS}}
\newcommand{\sym}[1]{\begin{block}{变量与来源}\small #1\end{block}}
\newcommand{\takeaway}[1]{\begin{exampleblock}{本页结论}\small #1\end{exampleblock}}
\newcommand{\sourceline}[1]{\par{\scriptsize\color{gray}#1}\par}
\title[方案 A：YYPF 替代 Fake]{方案 A：用随机 \ensuremath{P_{\rm ref}} proposal 替代 Fake score}
\subtitle{RTDMD 少步蒸馏中的 YYPF 条件均值估计与路径律校正}
\author{组会讲解版}
\date{2026 年 10 月}
\hypersetup{pdftitle={RTDMD Scheme A reorganized},pdfauthor={YYPF}}


% Standalone vector illustration of gradual denoising; no external image files.
\newcommand{\noisepic}[1]{%
\begin{tikzpicture}[x=0.105cm,y=0.105cm]
\fill[blue!15] (0,0) rectangle (20,20);
\fill[green!40!black] (0,0) rectangle (20,5);
\fill[yellow!75!orange] (5,5) rectangle (15,12);
\fill[red!65!black] (3,12)--(10,18)--(17,12)--cycle;
\fill[brown!80!black] (9,5) rectangle (12,10);
\fill[blue!55] (6,8) rectangle (8,10);
\pgfmathsetseed{51}
\foreach \xx in {0,...,19}{\foreach \yy in {0,...,19}{%
\pgfmathtruncatemacro{\shade}{100*rnd}
\fill[black!\shade,opacity=#1] (\xx,\yy) rectangle ++(1,1);}}
\draw[gray] (0,0) rectangle (20,20);
\end{tikzpicture}}

\begin{document}
\maketitle



\section{背景：为什么需要替代 Fake}

\begin{frame}{RTDMD 在一轮训练里做什么？}
\small
\begin{columns}[T,onlytextwidth]
\column{.54\textwidth}
\begin{enumerate}\itemsep3pt
\item 当前 Student $\theta_n$ 先给出一个中间状态 $X_{\theta_n}$。
\item 采样时间 $t$ 和查询噪声 $\varepsilon_q$，形成
\begin{equation}\label{eq:query-definition}
 x_t=\alpha_tX_{\theta_n}+\sigma_t\varepsilon_q.
\end{equation}
\item \ensuremath{P_{\rm ref}} 和 Student 都在同一个 $(x_t,t,c)$ 上评价。
\item 计算 score 差，并用 surrogate gradient 更新 Student。
\end{enumerate}
\begin{block}{RTDMD 型更新}
若 Student 输出 clean prediction $X_{{\rm pred},\theta_n}$，则可写成
{\small
\begin{equation}
\begin{aligned}
 g_n&=\lambda_t (J_{\theta_n}X_{{\rm pred},\theta_n})^{\top}\\
 &\quad\cdot
 \operatorname{sg}\!\left[\widehat s_{\theta_n}(x_t,t)-s_\mathrm{ref}(x_t,t)\right],\\[-1mm]
 \theta_{n+1}&=\theta_n-\eta_n g_n.
\end{aligned}
\end{equation}
}
\end{block}
\column{.44\textwidth}
\begin{tikzpicture}[node distance=3mm, every node/.style={font=\scriptsize}]
\node[draw=deepblue,fill=softblue,rounded corners,align=center,inner sep=3pt,text width=31mm] (s) {Student $\theta_n$\\中间状态 $X_{\theta_n}$};
\node[draw=deepblue,fill=softblue,rounded corners,align=center,inner sep=3pt,text width=31mm,below=of s] (q) {查询 $x_t=\alpha_tX_{\theta_n}+\sigma_t\varepsilon_q$};
\node[draw=deepgreen,fill=softgreen,rounded corners,align=center,inner sep=3pt,text width=31mm,below=of q] (ys) {Student 条件 score\\原 fake 或 YYPF};
\node[draw=deeporange,fill=softorange,rounded corners,align=center,inner sep=3pt,text width=31mm,below=of ys] (ts) {\ensuremath{P_{\rm ref}} score\\$s_\mathrm{ref}(x_t,t,c)$};
\node[draw=purple,fill=white,rounded corners,align=center,inner sep=3pt,text width=31mm,below=of ts] (u) {score 差\\更新 $\theta$};
\draw[-Latex,thick] (s)--(q);
\draw[-Latex,thick] (q)--(ys);
\draw[-Latex,thick] (ys)--(ts);
\draw[-Latex,thick] (ts)--(u);
\end{tikzpicture}
\end{columns}
\end{frame}

\begin{frame}{方案 A 要替代什么}
\begin{block}{RTDMD 需要的 Student score}
给定同一个 $(x_t,t,c)$，RTDMD 需要 Student 的条件 score
\begin{equation}\label{eq:student-target}
 s_\theta(x_t,t)=\nabla_{x_t}\log p_{\theta,t}(x_t)
 =\frac{\alpha_t m_{\theta,t}(x_t)-x_t}{\sigma_t^2},
 \qquad m_{\theta,t}(x_t)=\E_\theta[X_{\rm src}\mid X_t=x_t].
\end{equation}
\end{block}
\begin{columns}[T,onlytextwidth]
\column{.46\textwidth}
\begin{alertblock}{原来的 fake}
用一个额外的 fake 网络或 fake score 直接近似 $m_{\theta,t}$ 或 $s_\theta$。
\end{alertblock}
\column{.08\textwidth}
\centering\Large $\longrightarrow$
\column{.46\textwidth}
\begin{exampleblock}{方案 A 的替代}
用 \ensuremath{P_{\rm ref}} 候选路径及 Student/\ensuremath{P_{\rm ref}} 路径比估计 $m_{\theta,t}$，再用 Tweedie 恒等式恢复 score；\ensuremath{P_{\rm ref}} 提供随机 proposal。
\end{exampleblock}
\end{columns}
{\small $X_{\rm src}$ 是 Student 的源 latent，$x_t$ 是其加噪查询点；固定条件 $c$。下文用干净端点讲解；若源 latent 是中间状态，应在对应节点停止。}\par

\end{frame}

\begin{frame}{原 Fake 的完整循环，和方案 A 的替代位置}
\small
\begin{center}
\begin{tikzpicture}[node distance=4mm,every node/.style={font=\small},
b/.style={draw=deepblue,fill=softblue,rounded corners,align=center,text width=39mm,inner sep=4pt}]
\node[b] (a) {当前 Student 样本\\构造加噪查询};
\node[b,right=of a] (b) {训练 Fake\\DSM 加一致性正则};
\node[b,right=of b] (c) {score gap 更新 Student\\分布变化后继续追踪};
\draw[-Latex,thick] (a)--(b);\draw[-Latex,thick] (b)--(c);
\end{tikzpicture}
\end{center}
\begin{equation}\label{eq:fake-total}
\mathcal L_{\rm fake,total}=\mathcal L_{\rm DSM}+\gamma\mathcal L_{\rm cons},
\qquad\gamma\geq0.
\end{equation}
\begin{columns}[T,onlytextwidth]
\column{.48\textwidth}
\begin{block}{原 Fake：学习一个随时可调用的估计器}
反复训练参数 $\phi$，用一次网络前向近似式 \eqref{eq:fake-optimum} 的条件平均。
Student 变化后，Fake 需要继续追踪。
\end{block}
\column{.48\textwidth}
\begin{exampleblock}{方案 A：在本次查询上估计条件平均}
\ensuremath{P_{\rm ref}} 提供候选路径，正确的路径权重提供 Student 后验均值，再代入同一个 score 公式。
\end{exampleblock}
\end{columns}
\begin{equation}\label{eq:replacement-interface}
s_\phi(x_t,t)\quad\longrightarrow\quad
\widehat m_{\theta,t}(x_t)\quad\longrightarrow\quad
\widehat s_\theta(x_t,t)=\frac{\alpha_t\widehat m_{\theta,t}(x_t)-x_t}{\sigma_t^2}.
\end{equation}
\sourceline{$\mathcal L_{\rm cons}$ 稳定 Fake 训练；替换 Fake 后无需把它另加到 Student。}
\end{frame}

\begin{frame}{Fake 怎样训练：用已知标签学习未知的边缘 score}
\small
Student 产生源状态 $X_{\rm src}$，按 $X_t=\alpha_tX_{\rm src}+\sigma_t\varepsilon_q$ 独立加噪。高斯核给出可直接计算的标签：
\begin{equation}\label{eq:dsm-label}
q_t(x\mid x_{\rm src})=\N(x;\alpha_tx_{\rm src},\sigma_t^2I),\qquad
\nabla_x\log q_t(x\mid x_{\rm src})=\frac{\alpha_tx_{\rm src}-x}{\sigma_t^2}.
\end{equation}
\begin{block}{只更新 Fake 的参数 $\phi$}
\begin{equation}\label{eq:fake-dsm}
\mathcal L_{\rm DSM}(\phi)=
\E\!\left[\omega(t)\left\|s_\phi(X_t,t)-
\frac{\alpha_tX_{\rm src}-X_t}{\sigma_t^2}\right\|^2\right],\qquad\omega(t)>0.
\end{equation}
训练数据来自当前 Student；在这一步把 Student 样本固定，只拟合 $s_\phi$。
\end{block}
\begin{exampleblock}{平方回归为什么能学到 Student score}
\begin{equation}\label{eq:fake-optimum}
s_\phi^*(x_t,t)=\E\!\left[\frac{\alpha_tX_{\rm src}-x_t}{\sigma_t^2}\middle|X_t=x_t\right]
=\frac{\alpha_tm_{\theta,t}(x_t)-x_t}{\sigma_t^2}=s_\theta(x_t,t).
\end{equation}
$m_{\theta,t}(x_t)$ 就是给定查询点后，源 latent $X_{\rm src}$ 的条件均值。Fake 学的是条件平均，并非某一个样本的标签。
\end{exampleblock}
\sourceline{依据：RTDMD 附录 D。}
\end{frame}

\begin{frame}{核心思想 1：先明确 Student 要估计的对象}
\small
\textbf{查询点固定为} $x_t$（以及条件 $c$）。Student 的随机变量统一记为 $X$。
式 \eqref{eq:student-target} 定义的目标，就是 Student 在查询点 $x_t$ 下的源端点条件均值。

\begin{equation}
\boxed{
m_{\theta,t}(x_t)
=
\E_{P_\theta}[X_{\rm src}\mid X_t=x_t]
=
\E_{P_\theta}[X_0\mid X_t=x_t]
=
\int x_0\,p_\theta(x_0\mid x_t)\,\dd x_0 .
}
\label{eq:conditional-mean}
\end{equation}

\begin{block}{式 \eqref{eq:student-target} 与式 \eqref{eq:conditional-mean} 的关系}
式 \eqref{eq:student-target} 与式 \eqref{eq:conditional-mean} 是\textbf{同一个目标}。式 \eqref{eq:conditional-mean} 只是把式 \eqref{eq:student-target} 中的
$X_{\rm src}$ 进一步记为 Student 的源端点 $X_0$，并把条件期望写成密度积分。
因此右边出现的变量应是 Student 的 $X_0$，而不是 $Y_0$。
\end{block}

为了用 $P_{\rm ref}$ 产生样本，接下来把积分变量改成 $P_{\rm ref}$ rollout
产生的端点 $Y_0$。这是\textbf{同一个状态空间上的变量替换}，不是把 Student
变量改名成 $P_{\rm ref}$ 变量。
\end{frame}

\begin{frame}{核心思想 2：由式 \eqref{eq:conditional-mean} 得到端点重要性采样}
\small
从式 \eqref{eq:conditional-mean} 的积分形式出发。用 $P_{\rm ref}$ 的端点条件密度作为 proposal，
在积分中乘上
$p_{\rm ref}(y_0\mid x_t)/p_{\rm ref}(y_0\mid x_t)$，得到
\begin{equation}
\boxed{
\begin{aligned}
m_{\theta,t}(x_t)
&=
\int
y_0\,
\frac{p_\theta(y_0\mid x_t)}
     {p_{\rm ref}(y_0\mid x_t)}
p_{\rm ref}(y_0\mid x_t)\,\dd y_0  \\
&=
\E_{P_{\rm ref}}
\left[
\frac{p_\theta(Y_0\mid x_t)}
     {p_{\rm ref}(Y_0\mid x_t)}
Y_0
\,\middle|\,Y_t=x_t
\right].
\end{aligned}
}
\label{eq:endpoint-is}
\end{equation}

因此定义 endpoint weight：
\begin{equation}
\boxed{
w_{\rm end}(Y_0;x_t)
=
\frac{p_\theta(Y_0\mid x_t)}
     {p_{\rm ref}(Y_0\mid x_t)} .
}
\label{eq:endpoint-weight}
\end{equation}

\begin{block}{符号约定}
Student 的随机变量始终写 $X$；$P_{\rm ref}$ rollout 的随机变量始终写 $Y$。
所以目标是 $\E_{P_\theta}[X_0\mid X_t=x_t]$，而实际加权的是 $P_{\rm ref}$
产生的 $Y_0$。$p_\theta(Y_0\mid x_t)$ 只是 Student 密度在该样本位置的取值。
\end{block}
\end{frame}

\begin{frame}{核心思想 3：为什么不能直接算 endpoint weight？}
\small
式 \eqref{eq:endpoint-weight} 所需的是
$w_{\rm end}(Y_0;x_t)=p_\theta(Y_0\mid x_t)/p_{\rm ref}(Y_0\mid x_t)$。
但 Student 与 $P_{\rm ref}$ 的最终 endpoint marginal density 通常不能直接计算。

\begin{block}{方案 A 的核心操作}
\textbf{不要直接计算 endpoint ratio。}
让 $P_{\rm ref}$ 从查询点 $x_t$ 出发生成一整条去噪路径，
然后计算 Student 与 $P_{\rm ref}$ 在整条路径上的 likelihood ratio。
\end{block}

固定 $P_{\rm ref}$ 的起点为 $Y_K=x_t$，因此一条 rollout 的状态顺序为
$Y_K\to Y_{K-1}\to\cdots\to Y_1\to Y_0$。

\begin{block}{为什么这样做可行？}
完整路径的概率可以按每一步 transition kernel 相乘；
因此 path ratio 可以由每一步的 kernel ratio 直接计算，而不需要先知道
最终的 endpoint marginal。
\end{block}

\textbf{下一步：}由 Markov 分解把 endpoint ratio 改写成可计算的 path ratio。
\end{frame}

\begin{frame}{核心思想 4：由式 \eqref{eq:endpoint-is}--\eqref{eq:endpoint-weight} 转到完整路径权重}
\footnotesize
固定 $Y_K=x_t$，令中间路径为 $H=(Y_1,\ldots,Y_{K-1})$。
由 Markov 分解，
\begin{equation}
\boxed{
p_\theta(y_{0:K-1}\mid x_t)
=
\prod_{k=1}^{K}M_k^\theta(y_{k-1}\mid y_k),
\qquad
p_{\rm ref}(y_{0:K-1}\mid x_t)
=
\prod_{k=1}^{K}M_k^{\rm ref}(y_{k-1}\mid y_k).
}
\label{eq:path-factorization}
\end{equation}

因此由式 \eqref{eq:path-factorization}，一条 $P_{\rm ref}$ rollout 的 path weight 为
\begin{equation}
\boxed{
\begin{aligned}
\Lambda_A
&=
\frac{p_\theta(Y_{0:K-1}\mid x_t)}
     {p_{\rm ref}(Y_{0:K-1}\mid x_t)}
=
\prod_{k=1}^{K}
\frac{M_k^\theta(Y_{k-1}\mid Y_k)}
     {M_k^{\rm ref}(Y_{k-1}\mid Y_k)}
\\
&=
\frac{p_\theta(Y_0\mid x_t)}
     {p_{\rm ref}(Y_0\mid x_t)}
\frac{p_\theta(H\mid Y_0,x_t)}
     {p_{\rm ref}(H\mid Y_0,x_t)} .
\end{aligned}
}
\label{eq:path-weight}
\end{equation}

第二行只是把完整路径写成 $(Y_0,H)$，即“端点 $\times$ 中间 bridge”。
因此式 \eqref{eq:path-weight} 已经把式 \eqref{eq:endpoint-weight} 的 endpoint weight
与一个中间路径的 bridge ratio 联系起来。

\begin{block}{下一步}
对固定 $Y_0=y_0$ 和 $Y_t=x_t$，在 $P_{\rm ref}$ 的 bridge distribution
下对式 \eqref{eq:path-weight} 取条件期望，就能把中间路径积分掉。
\end{block}
\end{frame}

\begin{frame}{核心思想 5：由 path weight 边缘化中间路径，恢复 endpoint ratio}
\small
固定 $Y_0=y_0$ 和 $Y_t=x_t$，在 $P_{\rm ref}$ 的 bridge distribution
$p_{\rm ref}(H\mid y_0,x_t)$ 下，对式 \eqref{eq:path-weight} 取条件期望：
\begin{equation}
\boxed{
\begin{aligned}
\E_{P_{\rm ref}}[\Lambda_A\mid Y_0=y_0,Y_t=x_t]
&=
\frac{p_\theta(y_0\mid x_t)}
     {p_{\rm ref}(y_0\mid x_t)}
\E_{P_{\rm ref}}
\left[
\frac{p_\theta(H\mid y_0,x_t)}
     {p_{\rm ref}(H\mid y_0,x_t)}
\middle|Y_0=y_0,Y_t=x_t
\right] \\
&=
\frac{p_\theta(y_0\mid x_t)}
     {p_{\rm ref}(y_0\mid x_t)}
\int
p_\theta(H\mid y_0,x_t)\,\dd H \\
&=
\frac{p_\theta(y_0\mid x_t)}
     {p_{\rm ref}(y_0\mid x_t)} .
\end{aligned}
}
\label{eq:endpoint-recovery}
\end{equation}

其中第二步使用了
$H\sim p_{\rm ref}(H\mid y_0,x_t)$，所以 $P_{\rm ref}$ density 在积分中消去，
最后利用 Student 的条件路径密度归一化为 $1$。

\begin{block}{式 \eqref{eq:endpoint-recovery} 的结论}
\textbf{path weight 的 bridge 条件平均 = endpoint weight。}
因此式 \eqref{eq:endpoint-weight} 中难以直接计算的 endpoint ratio，可以通过整条
$P_{\rm ref}$ rollout 的 path ratio 来实现。
\end{block}
\end{frame}

\begin{frame}{核心思想 6：由 endpoint ratio 提升到有限 rollout 的路径估计}
\small
把式 \eqref{eq:endpoint-recovery} 代回式 \eqref{eq:endpoint-is}，等价地把 endpoint importance sampling 提升到路径空间：
\begin{equation}
\boxed{
m_{\theta,t}(x_t)
=
\frac{
\E_{P_{\rm ref}}[\Lambda_A Y_0\mid Y_t=x_t]
}{
\E_{P_{\rm ref}}[\Lambda_A\mid Y_t=x_t]
}.
}
\label{eq:path-identity}
\end{equation}

实际只生成 $N$ 条 $P_{\rm ref}$ rollout。对第 $i$ 条路径，
$Y_K^{(i)}=x_t$，并逐步采样
$Y_{k-1}^{(i)}\sim M_k^{\rm ref}(\cdot\mid Y_k^{(i)})$。
根据式 \eqref{eq:path-weight}，其 log-weight 为
$\ell^{(i)}=\sum_k\log\!\bigl[M_k^\theta(Y_{k-1}^{(i)}\mid Y_k^{(i)})/
M_k^{\rm ref}(Y_{k-1}^{(i)}\mid Y_k^{(i)})\bigr]$。

对有限 rollout 使用 self-normalized Monte Carlo，得到
\begin{equation}
\boxed{
\widehat m_{\theta,t}(x_t)
=
\frac{\sum_{i=1}^{N}\Lambda_A^{(i)}Y_0^{(i)}}
     {\sum_{i=1}^{N}\Lambda_A^{(i)}}
=
\sum_{i=1}^{N}\widetilde w_iY_0^{(i)},
\qquad
\widetilde w_i
=
\frac{\Lambda_A^{(i)}}{\sum_{j=1}^{N}\Lambda_A^{(j)}} .
}
\label{eq:mc-estimator}
\end{equation}

\begin{exampleblock}{实际计算链}
$P_{\rm ref}$ 负责 rollout $\rightarrow$ 式 \eqref{eq:path-weight} 计算 path weight
$\rightarrow$ 式 \eqref{eq:endpoint-recovery} 说明它对应 endpoint weight $\rightarrow$
式 \eqref{eq:mc-estimator} 得到 Student 条件均值估计。
\end{exampleblock}
\end{frame}






\begin{frame}{方案 A：进入实际计算}
\begin{center}
\vfill
{\Large\bfseries 从理论恒等式进入可计算的路径权重}
\vspace{5mm}
\begin{block}{接下来计算什么？}
固定查询点 $x_t$，由 $P_{\rm ref}$ 生成候选路径；随后逐步计算
$M_k^\theta(Y_{k-1}\mid Y_k)/M_k^{\rm ref}(Y_{k-1}\mid Y_k)$，
并累乘得到 path weight $\Lambda_A$，再用归一化权重估计 Student 条件均值。
\end{block}
\vfill
\end{center}
\end{frame}
\begin{frame}{步骤 1：得到 \ensuremath{P_{\rm ref}} 的候选去噪路径}
\small
\ensuremath{P_{\rm ref}} 从带噪图像 $Y_t$ 出发，逐步降低噪声，最终到达干净图像 $Y_0$。以四步采样为例：
\begin{center}
\begin{tikzpicture}[node distance=5mm,every node/.style={font=\small}]
\node (a) {\noisepic{0.94}};
\node[right=of a] (b) {\noisepic{0.70}};
\node[right=of b] (c) {\noisepic{0.43}};
\node[right=of c] (d) {\noisepic{0.18}};
\node[right=of d] (e) {\noisepic{0}};
\node[above=1mm of a] {$Y_t=Y_4$};\node[above=1mm of b] {$Y_3$};
\node[above=1mm of c] {$Y_2$};\node[above=1mm of d] {$Y_1$};\node[above=1mm of e] {$Y_0$};
\draw[-Latex,thick] (a)--(b);\draw[-Latex,thick] (b)--(c);
\draw[-Latex,thick] (c)--(d);\draw[-Latex,thick] (d)--(e);
\node[below=0mm of a] {带噪起点};\node[below=0mm of e] {干净终点};
\end{tikzpicture}
\end{center}
\vspace{2mm}
一般用 $K$ 表示步数，状态下标按 $K,K-1,\ldots,0$ 递减：
\begin{equation}\label{eq:ref-path}
Y_{k-1}^{(i)}\sim M_k^\mathrm{ref}(\cdot\mid Y_k^{(i)}),\qquad
Y^{(i)}=(Y_K^{(i)},\ldots,Y_0^{(i)}),\quad i=1,\ldots,N.
\end{equation}
保存每一步的起点、终点和 \ensuremath{P_{\rm ref}} 预测均值。Student 随后在这些路径点上评价。

\end{frame}

\begin{frame}{步骤2：求 YYPF 权重 $w$\\[-1mm]\normalsize 步骤2.1：计算 Student 的漂移项}
\footnotesize
将 \ensuremath{P_{\rm ref}} 路径点 $Y_k$ 输入 Student 网络 $D_\theta$。$D_\theta(Y_k,s_k)$ 是对干净端点的预测（clean prediction）；由它重构旧噪声估计：
\begin{equation}\label{eq:cps-prediction}
X_{{\rm pred},k}=D_\theta(Y_k,s_k)=\widehat X_0^{(k)},\qquad
\widehat\varepsilon_k^{\rm old}=
\frac{Y_k-(1-s_k)X_{{\rm pred},k}}{s_k}.
\end{equation}
沿用式 \eqref{eq:cps-prediction} 的记号，将式 \eqref{eq:cps-prediction} 中的 $X_{{\rm pred},k}$ 暂记为 $\widehat X$。\par
将旧噪声的定义乘以 $s_k$，移项后得到：
\begin{equation}\label{eq:state-decomposition}
\boxed{Y_k=(1-s_k)\widehat X+s_k\widehat\varepsilon_k^{\rm old}.}
\end{equation}
\begin{block}{节点顺序与时间方向}
\begin{equation}
s_K=t>s_{K-1}>\cdots>s_0=0,\qquad
\Delta\tau_k=s_k-s_{k-1}>0.
\end{equation}
$s_k$ 是当前点的噪声水平，$s_{k-1}$ 是下一点的噪声水平；$\tau$ 随反向采样递增。每一步从 $k$ 到 $k-1$；噪声水平下降，采样时钟 $\tau$ 增加。
\end{block}
{\small \ensuremath{P_{\rm ref}} 状态用 $Y$，Student 状态用 $X$；参数仍用 $\mathrm{ref},\theta$。$D_\theta(Y_k,s_k)$ 是 Student 在 \ensuremath{P_{\rm ref}} 状态上的评价。}
\end{frame}


\begin{frame}{步骤2.1：计算 Student 的漂移项}
\begin{center}
\vfill
{\Large\bfseries 目标：得到 Student 每一步的有效漂移}
\vspace{5mm}
\begin{block}{本阶段输出}
从 Student 网络在 $P_{\rm ref}$ 路径点上的预测出发，得到一步均值
$\mu_{\theta,k}(Y_k)$，进一步得到 Student 有效漂移 $b_{\theta,k}(Y_k)$
以及可评价的 Student 转移核 $M_k^\theta$。
\end{block}
\vfill
\end{center}
\end{frame}
\begin{frame}{步骤2.1：计算 Student 的漂移项}
\small
在当前点令 $X_k=Y_k$。按线性时间表，下一节点的信号系数是 $1-s_{k-1}$，噪声系数是 $s_{k-1}$：
\begin{equation}\label{eq:deterministic-cps}
X_{k-1}=(1-s_{k-1})\widehat X+s_{k-1}\widehat\varepsilon_k^{\rm old}.
\end{equation}
代入式 \eqref{eq:cps-prediction} 的旧噪声定义并合并 $\widehat X$ 的系数：
\begin{equation}\label{eq:deterministic-interpolation}
\begin{aligned}
X_{k-1}
&=(1-s_{k-1})\widehat X+\frac{s_{k-1}}{s_k}\bigl[Y_k-(1-s_k)\widehat X\bigr]\\
&=\boxed{\frac{s_{k-1}}{s_k}Y_k+\left(1-\frac{s_{k-1}}{s_k}\right)\widehat X}.
\end{aligned}
\end{equation}
\begin{block}{从当前位置向预测端点移动一段距离}
$X_{k-1}$ 表示 Student 在噪声水平 $s_{k-1}$ 的下一步状态，也可记作 $X_{s_{k-1}}$。这里没有新噪声，给定 $Y_k$ 后结果确定；加入 CPS 随机噪声后，它表示实际抽到的下一步状态。
\end{block}
由于 $0\le s_{k-1}/s_k<1$，这一步位于 $Y_k$ 与预测端点 $\widehat X$ 之间。接下来保留时间表，改用新旧噪声的混合来生成不同路径。
\end{frame}

\begin{frame}{步骤2.1：计算 Student 的漂移项}
\small
CPS（coefficient-preserving sampling，系数保持采样）统一了 Euler 与一致性模型（CM）的多步采样，沿用 CM 的“预测信号、再加噪”方式。

为了让同一当前状态产生不同路径，先将旧噪声替换为：
\begin{equation}\label{eq:cps-noise-mix}
\zeta_k=\cos(\eta\pi/2)\widehat\varepsilon_k^{\rm old}
+\sin(\eta\pi/2)\xi_k,\qquad\xi_k\sim\N(0,I).
\end{equation}
再按下一节点的系数组合信号与混合噪声：
\begin{equation}\label{eq:cps-combination}
X_{k-1}=(1-s_{k-1})\widehat X+s_{k-1}\zeta_k.
\end{equation}
代入式 \eqref{eq:cps-noise-mix}，并用 $X_{{\rm pred},k}$ 写回 $\widehat X$：
\begin{equation}\label{eq:cps-step}
\boxed{X_{k-1}=(1-s_{k-1})X_{{\rm pred},k}
+s_{k-1}\cos(\eta\pi/2)\widehat\varepsilon_k^{\rm old}
+s_{k-1}\sin(\eta\pi/2)\xi_k.}
\end{equation}
$\eta\in[0,1]$ 控制新噪声比例；$\xi_k$ 与当前状态和旧噪声估计独立。$X_{k-1}$ 是加上本步随机噪声后，实际得到的下一步状态。
\sourceline{RTDMD 附录 C；CPS 并不要求 Student 本身是一致性模型。}
\end{frame}

\begin{frame}{步骤2.1：计算 Student 的漂移项}
\small
给定 $X_k=Y_k$，端点预测和旧噪声估计固定，新噪声均值为零。因此 CPS 的下一步平均位置为：
\begin{equation}\label{eq:cps-mean}
\mu_{\theta,k}(Y_k)=\E[X_{k-1}\mid X_k=Y_k]
=(1-s_{k-1})X_{{\rm pred},k}
+s_{k-1}\cos(\eta\pi/2)\widehat\varepsilon_k^{\rm old}.
\end{equation}
$\mu_{\theta,k}$ 是 \textbf{CPS 下一步的平均位置}，不是网络的干净端点预测。记新噪声的系数为 $\nu_k$，则
\begin{equation}\label{eq:cps-sample-mean}
\nu_k=s_{k-1}\sin(\eta\pi/2),\qquad
X_{k-1}=\mu_{\theta,k}(Y_k)+\nu_k\xi_k.
\end{equation}
\begin{equation}\label{eq:cps-kernel}
M_k^\theta(x'\mid X_k=Y_k)=\N(x';\mu_{\theta,k}(Y_k),\nu_k^2I).
\end{equation}
\begin{center}
\begin{tabular}{ll}
\toprule
$D_\theta(Y_k,s_k)$ & 网络预测的干净端点\\
$\mu_{\theta,k}(Y_k)$ & CPS 下一步的平均位置\\
$X_{k-1}$ & 加入本步随机噪声后实际得到的下一步状态\\
\bottomrule
\end{tabular}
\end{center}
\sourceline{这里固定当前参数 $\theta$，省略训练轮次 $n$；$X_k$ 是 $X_{s_k}$ 的简写。$s_{k-1}=0$ 时末步确定，不适用非退化高斯密度式。}
\end{frame}

\begin{frame}{步骤2.1：计算 Student 的漂移项}
\small
上一步已经得到 $\mu_{\theta,k}(Y_k)$ 和 $\nu_k^2I$。接下来要与 \ensuremath{P_{\rm ref}} 比较漂移，先把 Student 的平均位移写出来。
\begin{equation}\label{eq:cps-displacement}
\E[X_{k-1}-Y_k\mid X_k=Y_k]=\mu_{\theta,k}(Y_k)-Y_k.
\end{equation}
一步耗时为 $\Delta\tau_k=s_k-s_{k-1}>0$。平均位移除以耗时，就是该步的有效漂移：
\begin{equation}\label{eq:cps-drift}
b_{\theta,k}^{\rm CPS}(Y_k)=\frac{\mu_{\theta,k}(Y_k)-Y_k}{\Delta\tau_k},\qquad
G_k=\frac{\nu_k}{\sqrt{\Delta\tau_k}}I.
\end{equation}
\begin{block}{现在可以用同一种形式比较 \ensuremath{P_{\rm ref}} 和 Student}
\begin{equation}\label{eq:student-euler}
X_{k-1}=\underbrace{Y_k+b_{\theta,k}^{\rm CPS}(Y_k)\Delta\tau_k}_{\text{Student 一步均值}}
+\underbrace{G_k\sqrt{\Delta\tau_k}\,\xi_k}_{\text{一步随机波动}}.
\end{equation}
这只是 CPS 离散核的等价改写，尚未假设它存在连续时间极限。
\end{block}
下一页将这两个表达逐项对照，并说明它与 SDE 离散化的关系。
\end{frame}

\begin{frame}{步骤2.1：计算 Student 的漂移项}
\small
为了使用 Girsanov 路径赋权，\ensuremath{P_{\rm ref}} proposal 与 Student target 都采用随机动力学：
\begin{align}
\dd \bar Y_\tau&=b_{\rm ref}(\tau,\bar Y_\tau)\dd\tau+G_\tau\dd W_\tau^\mathrm{ref}
&&\text{proposal: }P_{\rm ref},\label{eq:ref-sde-overview}\\
\dd \bar X_\tau&=b_\theta(\tau,\bar X_\tau)\dd\tau+G_\tau\dd W_\tau^\theta
&&\text{Student：}P_\theta,\label{eq:student-sde-overview}
\end{align}
$\tau$ 是递增的采样时钟，$\tau_k=t-s_k$；$\bar Y_{\tau_k}=Y_k$、$\bar X_{\tau_k}=X_k$。$W^\mathrm{ref},W^\theta$ 分别是各自路径律下的布朗运动。
\begin{block}{两边都加噪声，并采用相同的扩散系数}
Student 使用 CPS 随机转移，\ensuremath{P_{\rm ref}} 也需指定随机 proposal。这里让两者采用同一个 $G_\tau$，使单位时间噪声协方差都是 $G_\tau G_\tau^\top$，只比较漂移差。

“相同”指噪声强度与协方差相同，不要求两者使用同一次随机取值。实际只采样 \ensuremath{P_{\rm ref}} 路径，Student 在这条路径上评价转移密度。
\end{block}
\small 原 \ensuremath{P_{\rm ref}} ODE 不会自动满足这两式；本方案另外构造随机 \ensuremath{P_{\rm ref}}。对于实际 CPS，先按共同网格和随机强度匹配单步协方差；能否取连续极限还需额外条件。下一页先给出 SDE 的离散形式。
\end{frame}

\begin{frame}{步骤2.1：计算 Student 的漂移项}
\small
对一般 Student SDE（式 \eqref{eq:student-sde}），在一个小区间内冻结漂移与扩散系数。Euler--Maruyama 离散给出：
\begin{equation}\label{eq:euler-general}
X_{k-1}\approx X_k+b_\theta(\tau_k,X_k)\Delta\tau_k
+G_k\sqrt{\Delta\tau_k}\,\xi_k,\qquad \xi_k\sim\N(0,I).
\end{equation}
令 $X_k=Y_k$，与 CPS 式 \eqref{eq:cps-step} 对照：
\begin{equation}
\begin{aligned}
\text{CPS：}\quad X_{k-1}
&=\underbrace{(1-s_{k-1})X_{{\rm pred},k}+s_{k-1}\cos(\eta\pi/2)\widehat\varepsilon_k^{\rm old}}_{\mu_{\theta,k}(Y_k)}
+\underbrace{\nu_k\xi_k}_{\text{噪声}},\\[3mm]
\text{SDE 小步：}\quad X_{k-1}
&=\underbrace{Y_k+b_{\theta,k}(Y_k)\Delta\tau_k}_{\text{均值}}
+\underbrace{G_k\sqrt{\Delta\tau_k}\,\xi_k}_{\text{噪声}}.
\end{aligned}
\end{equation}
\begin{block}{匹配均值与协方差}
\begin{equation}
\mu_{\theta,k}=Y_k+b_{\theta,k}\Delta\tau_k,\qquad
G_kG_k^\top\Delta\tau_k=\nu_k^2I.
\end{equation}
因此式 \eqref{eq:cps-drift} 给出的漂移与扩散恰好还原 CPS 核。后面先用该核计算离散密度比，再解释它在连续时间下对应的 Girsanov 权重。
\end{block}
\end{frame}

\begin{frame}{步骤2.2：计算 \ensuremath{P_{\rm ref}} 的漂移项}
\small
\ensuremath{P_{\rm ref}} 原来的 ODE 是确定性的。要得到可计算的高斯 proposal，另外为 \ensuremath{P_{\rm ref}} 指定 CPS 随机采样规则。先把速度输出转成端点预测。
\begin{equation}\label{eq:ref-velocity}
Y_k=y,\qquad v_k=v_{\rm ref}(y,s_k).
\end{equation}
\begin{block}{采用 rectified flow 的速度约定}
\begin{equation}\label{eq:rf-convention}
y_s=(1-s)y_0+s\varepsilon,\qquad v=\varepsilon-y_0,
\qquad y_s=y_0+sv.
\end{equation}
因此，当前位置减去“噪声水平乘以速度”，得到 \ensuremath{P_{\rm ref}} 的干净端点预测：
\begin{equation}\label{eq:ref-prediction}
\boxed{y_{{\rm pred},k}^{\rm ref}=y-s_kv_k},\qquad
\widehat\varepsilon_{\mathrm{ref},k}^{\rm old}
=\frac{y-(1-s_k)y_{{\rm pred},k}^{\rm ref}}{s_k}.
\end{equation}
\end{block}
\ensuremath{P_{\rm ref}} 的端点预测用 $y_{{\rm pred},k}^{\rm ref}$ 表示，与 Student 的 $X_{{\rm pred},k}$ 区分。二者都是信号预测，随后分别送入相同形式的 CPS 更新。
\sourceline{这里的 $v_{\rm ref}$ 是随噪声水平 $s$ 变化的速度；它不等于采样时钟 $\tau$ 下的有效漂移 $b_{\mathrm{ref},k}$。}
\end{frame}

\begin{frame}{步骤2.2：计算 \ensuremath{P_{\rm ref}} 的漂移项}
\small
对 Student 的同一步 $s_k\to s_{k-1}$，选择相同的随机强度 $\eta>0$，并定义 \ensuremath{P_{\rm ref}} 的 CPS 均值：
\begin{equation}\label{eq:ref-cps-mean}
\mu_{\mathrm{ref},k}^{\rm prop}(y)=(1-s_{k-1})y_{{\rm pred},k}^{\rm ref}
+s_{k-1}\cos(\eta\pi/2)\widehat\varepsilon_{\mathrm{ref},k}^{\rm old},\qquad
\nu_k=s_{k-1}\sin(\eta\pi/2).
\end{equation}
\begin{equation}\label{eq:ref-cps-sample}
Y_{k-1}=\mu_{\mathrm{ref},k}^{\rm prop}(Y_k)+\nu_k\xi_k^\mathrm{ref},
\qquad \xi_k^\mathrm{ref}\sim\N(0,I).
\end{equation}
新噪声使这条采样规则具有高斯转移核，后文简记为 $M_k^\mathrm{ref}$：
\begin{equation}\label{eq:ref-cps-kernel}
\boxed{M_k^\mathrm{ref}(y'\mid y)\equiv M_{\mathrm{ref},k}^{\rm prop}(y'\mid y)
=\N(y';\mu_{\mathrm{ref},k}^{\rm prop}(y),\nu_k^2I)},\quad\nu_k>0.
\end{equation}
\begin{equation}\label{eq:ref-effective-drift}
b_{\mathrm{ref},k}(y)=\frac{\mu_{\mathrm{ref},k}^{\rm prop}(y)-y}{\Delta\tau_k},
\qquad G_kG_k^\top\Delta\tau_k=\nu_k^2I.
\end{equation}
\begin{block}{分子与分母使用同一个 Student 步长}
比值 $M_k^\theta(Y_{k-1}\mid Y_k)/M_k^\mathrm{ref}(Y_{k-1}\mid Y_k)$ 的两种核，都对应 Student 的第 $k$ 个区间 $s_k\to s_{k-1}$。\ensuremath{P_{\rm ref}} 也按这套 $K$ 步网格构造 proposal，不使用其原来多步 ODE 求解器的内部小步。
\end{block}
\sourceline{相同的 $\eta$ 与时间网格保证协方差相同，不要求两条路径使用相同的噪声实现。零噪声末步仍需单独处理。}
\end{frame}

\begin{frame}{步骤2.3：计算 \ensuremath{P_{\rm ref}} 到 Student 的重要性校正}
\footnotesize
\begin{block}{承接前面的核心思想：现在开始计算哪个量？}
由式 \eqref{eq:endpoint-recovery}，endpoint ratio 可由 path weight 的 bridge 条件平均恢复；式
\eqref{eq:path-weight} 给出 path weight 的定义。因此从本页开始，\textbf{进入实际计算阶段}：
我们实际计算每条 $P_{\rm ref}$ rollout 的 path weight $\Lambda_A$，再用它完成 Student 条件均值的路径空间估计。
\end{block}
\begin{block}{proposal 与 target}
\ensuremath{P_{\rm ref}} 产生的路径来自 $P_{\rm ref}$，但目标来自 Student $P_\theta$。重要性采样的基本修正是
\begin{equation}
\text{weight}=\frac{\text{目标密度}}{\text{提议密度}}.
\end{equation}
\end{block}
\begin{center}
\begin{tikzpicture}[node distance=8mm, every node/.style={font=\small},
 box/.style={draw=deepblue,fill=softblue,rounded corners,align=center,inner sep=5pt,text width=38mm}]
\node[box] (q) {\ensuremath{P_{\rm ref}} 产生\\候选路径 $P_{\rm ref}$};
\node[box,right=of q] (ratio) {乘上\\$P_\theta/P_{\rm ref}$};
\node[box,right=of ratio] (p) {加权后\\回到 Student $P_\theta$};
\draw[-Latex,thick] (q)--(ratio)--(p);
\end{tikzpicture}
\end{center}
\begin{exampleblock}{为什么不能直接平均}
直接平均只得到 \ensuremath{P_{\rm ref}} 的终点均值；ratio 会把 \ensuremath{P_{\rm ref}} 更常走的地方降权，把 Student 更偏好的地方升权。
\end{exampleblock}
\end{frame}

\begin{frame}{步骤2.3：计算 \ensuremath{P_{\rm ref}} 到 Student 的重要性校正}
\small
高斯密度中的残差是“实际下一步位置减去模型预测均值”。在同一个 \ensuremath{P_{\rm ref}} 小步上，分别用两种模型评价：
\begin{equation}\label{eq:residual-derivation}
\begin{aligned}
r_{a,k}&=Y_{k-1}-\bigl[Y_k+b_{a,k}(Y_k)\Delta\tau_k\bigr]\\
&=\underbrace{Y_{k-1}-Y_k}_{\Delta Y_k}-b_{a,k}(Y_k)\Delta\tau_k,
\qquad a\in\{\theta,\mathrm{ref}\}.
\end{aligned}
\end{equation}
$a=\theta$ 对应 Student，$a=\mathrm{ref}$ 对应 \ensuremath{P_{\rm ref}}；$\Delta Y_k=Y_{k-1}-Y_k$。两者评价同一个终点，但预测均值不同。
\begin{equation}\label{eq:gaussian-density}
\boxed{M_k^a(Y_{k-1}\mid Y_k)=
\frac{1}{(2\pi)^{d/2}|A_k\Delta\tau_k|^{1/2}}
\exp\!\left(-\frac12 r_{a,k}^{\top}(A_k\Delta\tau_k)^{-1}r_{a,k}\right).}
\end{equation}
\begin{block}{取对数}
\begin{equation}\label{eq:log-gaussian}
\log M_k^a=-\frac d2\log(2\pi)-\frac12\log|A_k\Delta\tau_k|
-\frac12r_{a,k}^{\top}(A_k\Delta\tau_k)^{-1}r_{a,k}.
\end{equation}
两种核协方差相同，作差时前两项抵消，得到下一页的式 \eqref{eq:one-step-ratio}。

$d$ 是状态维数，$A_k=G_kG_k^\top$，$|\cdot|$ 表示行列式；$a=\theta$ 对应 Student，$a=\mathrm{ref}$ 对应 \ensuremath{P_{\rm ref}}。
\end{block}
\end{frame}

\begin{frame}{步骤2.3：计算 \ensuremath{P_{\rm ref}} 到 Student 的重要性校正}
\small
将两种核的对数密度相减；共同协方差使归一化项抵消，留下残差二次型之差：

\begin{equation}\label{eq:one-step-ratio}
\log\rho_k\equiv\log\frac{M_k^\theta(Y_{k-1}\mid Y_k)}{M_k^\mathrm{ref}(Y_{k-1}\mid Y_k)}
=-\frac12\left[r_{\theta,k}^{\top}(A_k\Delta\tau_k)^{-1}r_{\theta,k}
-r_{\mathrm{ref},k}^{\top}(A_k\Delta\tau_k)^{-1}r_{\mathrm{ref},k}\right].
\end{equation}

一条路径包含 $K$ 个小步。将单步密度比相乘，得到整条路径的权重：
\begin{equation}
\Lambda_A^{(i)}=\prod_{k=1}^{K}
\frac{M_k^\theta(Y_{k-1}^{(i)}\mid Y_k^{(i)})}
{M_k^\mathrm{ref}(Y_{k-1}^{(i)}\mid Y_k^{(i)})}.
\end{equation}
相应的对数权重为：
\begin{equation}
\ell^{(i)}=\log\Lambda_A^{(i)}
=\sum_{k=1}^{K}\log\frac{M_k^\theta(Y_{k-1}^{(i)}\mid Y_k^{(i)})}
{M_k^\mathrm{ref}(Y_{k-1}^{(i)}\mid Y_k^{(i)})}.
\end{equation}

\small 上式要求每一步的密度比存在。CPS 的零噪声末步是确定映射；若两模型末步不同，不能直接套高斯核比，需要单独处理终端映射或采用可评价的条件路径核。
\end{frame}

\begin{frame}{步骤2.3：计算 \ensuremath{P_{\rm ref}} 到 Student 的重要性校正}
\small
离散核比已经能赋权。Girsanov 解释：在共同扩散下，为什么只调整漂移就能改变路径分布？
\begin{align}
\dd \bar Y_\tau&=b_{\rm ref}(\tau,\bar Y_\tau)\dd\tau+G_\tau\dd W_\tau^\mathrm{ref}
&&\text{proposal：}P_{\rm ref},\label{eq:ref-sde}\\
\dd \bar X_\tau&=b_\theta(\tau,\bar X_\tau)\dd\tau+G_\tau\dd W_\tau^\theta
&&\text{target：}P_\theta.\label{eq:student-sde}
\end{align}
这里 $\tau$ 是递增的采样时钟，$\tau_k=t-s_k$；$\bar Y_{\tau_k}=Y_k$、$\bar X_{\tau_k}=X_k$。实际模拟 \ensuremath{P_{\rm ref}}，在同一状态处评价两个漂移。
\begin{equation}\label{eq:u-definition}
\delta b_\tau=b_\theta(\tau,\bar Y_\tau)-b_{\rm ref}(\tau,\bar Y_\tau),\qquad
u_\tau=G_\tau^{-1}\delta b_\tau.
\end{equation}
\begin{block}{为什么把漂移差除以扩散矩阵}
$\delta b$ 表示 Student 比 \ensuremath{P_{\rm ref}} 每单位时间多移动多少；$u$ 用噪声尺度衡量这个差异。噪声很小时，同样的漂移差更难靠 proposal 的随机波动覆盖。
\end{block}
{\small 共同扩散是连续 Girsanov 的前提。本稿采用可逆 $G$，并要求绝对连续及适当可积性；例如 Novikov 条件是一个充分条件。}
\end{frame}

\begin{frame}{步骤2.3：计算 \ensuremath{P_{\rm ref}} 到 Student 的重要性校正}
\small
\ensuremath{P_{\rm ref}} 的实际移动由两部分组成：预测的平均移动，以及随机偏移。
\begin{equation}
\Delta Y_k=b_{\mathrm{ref},k}(Y_k)\Delta\tau_k+G_k\sqrt{\Delta\tau_k}\,\varepsilon_k^\mathrm{ref}.
\end{equation}
\begin{block}{先减去 \ensuremath{P_{\rm ref}} 已经预测到的部分}
\begin{equation}\label{eq:innovation-detail}
r_{\mathrm{ref},k}=\Delta Y_k-b_{\mathrm{ref},k}(Y_k)\Delta\tau_k
=G_k\sqrt{\Delta\tau_k}\,\varepsilon_k^\mathrm{ref}.
\end{equation}
\end{block}
\begin{equation}\label{eq:whiten}
\Delta\widehat W_k=G_k^{-1}r_{\mathrm{ref},k},\quad
\Cov(\Delta\widehat W_k\mid Y_k)=\Delta\tau_kI,\qquad
\widehat\varepsilon_k=\frac{\Delta\widehat W_k}{\sqrt{\Delta\tau_k}}\sim\N(0,I).
\end{equation}
{\small 权重应判断“这次随机偏移是否更符合 Student”。直接使用 $\Delta Y_k$ 会把 \ensuremath{P_{\rm ref}} 的平均移动也当作随机残差；标准正态噪声与 布朗增量相差 $\sqrt{\Delta\tau_k}$。}
\end{frame}

\begin{frame}{步骤2.3：计算 \ensuremath{P_{\rm ref}} 到 Student 的重要性校正}
\small
从式 \eqref{eq:one-step-ratio} 的高斯对数密度比出发，将 Student 残差写成 \ensuremath{P_{\rm ref}} 残差减去漂移差：
\begin{equation}
r_{\theta,k}=r_{\mathrm{ref},k}-\delta b_k \Delta\tau_k,\qquad
\delta b_k=b_{\theta,k}(Y_k)-b_{\mathrm{ref},k}(Y_k).
\end{equation}
\begin{equation}\label{eq:expanded-ratio}
\begin{aligned}
\log\rho_k
&=-\frac{(r_{\mathrm{ref},k}-\delta b_k\Delta\tau_k)^\top A_k^{-1}(r_{\mathrm{ref},k}-\delta b_k\Delta\tau_k)
-r_{\mathrm{ref},k}^\top A_k^{-1}r_{\mathrm{ref},k}}{2\Delta\tau_k}\\
&=\delta b_k^\top A_k^{-1}r_{\mathrm{ref},k}
-\frac12\delta b_k^\top A_k^{-1}\delta b_k \Delta\tau_k.
\end{aligned}
\end{equation}
\begin{block}{代入白化后的漂移差和 布朗增量}
\begin{equation}\label{eq:whitened-ratio}
\boxed{\log\rho_k=u_k^\top\Delta\widehat W_k-\tfrac12\|u_k\|^2\Delta\tau_k}
=\sqrt{\Delta\tau_k}\,u_k^\top\widehat\varepsilon_k-\tfrac12\|u_k\|^2\Delta\tau_k.
\end{equation}
\end{block}
式 \eqref{eq:expanded-ratio} 由式 \eqref{eq:one-step-ratio} 代入残差关系并展开二次型得到；接下来分别推导两项如何白化。
\end{frame}

\begin{frame}{步骤2.3：计算 \ensuremath{P_{\rm ref}} 到 Student 的重要性校正}
\small
式 \eqref{eq:expanded-ratio} 的第二项衡量漂移差的大小。先用噪声尺度统一度量：
\begin{equation}
A_k=G_kG_k^\top,\qquad
A_k^{-1}=G_k^{-\top}G_k^{-1},\qquad
u_k=G_k^{-1}\delta b_k.
\end{equation}
\begin{block}{把矩阵二次型写成向量长度的平方}
\begin{equation}
\begin{aligned}
\delta b_k^\top A_k^{-1}\delta b_k
&=\delta b_k^\top G_k^{-\top}G_k^{-1}\delta b_k\\
&=(G_k^{-1}\delta b_k)^\top(G_k^{-1}\delta b_k)
=u_k^\top u_k=\|u_k\|^2.
\end{aligned}
\end{equation}
\end{block}
因此，第二项直接变为
\begin{equation}
\boxed{-\frac12\delta b_k^\top A_k^{-1}\delta b_k\Delta\tau_k
=-\frac12\|u_k\|^2\Delta\tau_k.}
\end{equation}
$G_k$ 为可逆扩散矩阵；$A_k$ 是单位时间协方差；$\delta b_k$ 是同一 $Y_k$ 处的 Student 减 \ensuremath{P_{\rm ref}} 漂移；$u_k$ 是白化漂移差。
\end{frame}

\begin{frame}{步骤2.3：计算 \ensuremath{P_{\rm ref}} 到 Student 的重要性校正}
\small
式 \eqref{eq:expanded-ratio} 的第一项使用 \ensuremath{P_{\rm ref}} 实际小步中的随机残差：
\begin{equation}
\Delta Y_k=b_{\mathrm{ref},k}(Y_k)\Delta\tau_k
+G_k\sqrt{\Delta\tau_k}\,\varepsilon_k^\mathrm{ref},\qquad
r_{\mathrm{ref},k}=G_k\sqrt{\Delta\tau_k}\,\varepsilon_k^\mathrm{ref}.
\end{equation}
\begin{block}{代入残差，再用 $A_k^{-1}=G_k^{-\top}G_k^{-1}$}
\begin{equation}
\begin{aligned}
\delta b_k^\top A_k^{-1}r_{\mathrm{ref},k}
&=\delta b_k^\top G_k^{-\top}G_k^{-1}G_k
\sqrt{\Delta\tau_k}\,\varepsilon_k^\mathrm{ref}\\
&=(G_k^{-1}\delta b_k)^\top\sqrt{\Delta\tau_k}\,\varepsilon_k^\mathrm{ref}
=u_k^\top\Delta\widehat W_k.
\end{aligned}
\end{equation}
\end{block}
\begin{equation}
\Delta\widehat W_k=\sqrt{\Delta\tau_k}\,\varepsilon_k^\mathrm{ref}
=G_k^{-1}r_{\mathrm{ref},k},\qquad
\boxed{\log\rho_k=u_k^\top\Delta\widehat W_k
-\frac12\|u_k\|^2\Delta\tau_k.}
\end{equation}
$\varepsilon_k^\mathrm{ref}\sim\N(0,I)$ 是 \ensuremath{P_{\rm ref}} 的标准正态噪声；$\Delta\widehat W_k$ 是对应 布朗增量，协方差为 $\Delta\tau_k I$。第一项衡量实际随机偏移与漂移差的对齐程度。
\end{frame}

\begin{frame}{步骤2.3：计算 \ensuremath{P_{\rm ref}} 到 Student 的重要性校正}
\small
\begin{columns}[T,onlytextwidth]
\column{.48\textwidth}\begin{block}{随机对齐项}
\begin{equation}
u_k^\top\Delta\widehat W_k
\end{equation}
随机残差与漂移差同向时，Student 对这一步给出的相对密度更高，权重随之增大；反向时减小。
\end{block}
\column{.48\textwidth}\begin{exampleblock}{二次校正项}
\begin{equation}
-\tfrac12\|u_k\|^2\Delta\tau_k
\end{equation}
漂移差越大，需要扣除的量越大。它让指数权重具有正确的归一化，不能省略。
\end{exampleblock}\end{columns}
\begin{equation}
\E_{P_{\rm ref}}\!\left[\exp\!\left(u_k^\top\Delta\widehat W_k-\tfrac12\|u_k\|^2\Delta\tau_k\right)\middle|Y_k\right]=1.
\end{equation}
\begin{block}{一维例子：同样的 drift gap，不同随机移动得到不同权重}
\begin{equation}
\Delta\tau_k=0.1,\quad G_k=1,\quad b_{\theta,k}-b_{\mathrm{ref},k}=2:
\qquad\log\rho_k=2r_{\mathrm{ref},k}-0.2.
\end{equation}
残差为 $0.3$ 时，log-weight 为 $0.4$；残差为 $-0.3$ 时，log-weight 为 $-0.8$。Student 更支持第一种移动。
\end{block}
\end{frame}

\begin{frame}{步骤2.3：计算 \ensuremath{P_{\rm ref}} 到 Student 的重要性校正}
\small
沿每条实际 \ensuremath{P_{\rm ref}} 路径，累加刚才得到的单步 log-weight：
\begin{equation}\label{eq:path-girsanov-discrete}
\ell^{(i)}=\sum_{k=1}^{K}\left[(u_k^{(i)})^\top\Delta\widehat W_k^{(i)}
-\tfrac12\|u_k^{(i)}\|^2\Delta\tau_k\right],\qquad \Lambda_A^{(i)}=e^{\ell^{(i)}}.
\end{equation}
\begin{block}{当这些小步来自一个适定连续扩散的离散化}
\begin{equation}\label{eq:girsanov}
\Lambda_A=\frac{\mathrm dP_\theta}{\mathrm dP_{\rm ref}}
=\exp\!\left(\int_0^{\tau_*}u_\tau^\top\dd W_\tau^\mathrm{ref}
-\frac12\int_0^{\tau_*}\|u_\tau\|^2\dd\tau\right).
\end{equation}
\begin{equation}\label{eq:change-of-measure}
\E_{P_\theta}[f(\bar X_{\tau_*})]=\E_{P_{\rm ref}}[\Lambda_A f(\bar Y_{\tau_*})].
\end{equation}
\end{block}
{\small 这就是赋权的用途：用 \ensuremath{P_{\rm ref}} 样本估计目标路径下的平均。高斯核比与离散 Girsanov 写法计算同一份权重，不能重复相乘。}
\sourceline{固定 CPS 随机强度未必存在有限扩散的连续极限；CPS 实现优先使用其实际离散核，上述连续式提供满足条件时的理论解释。}
\end{frame}

\begin{frame}{步骤2.3：计算 \ensuremath{P_{\rm ref}} 到 Student 的重要性校正}
对 $N$ 条路径得到 $\ell^{(1)},\ldots,\ell^{(N)}$。先减去最大值 $\ell_{\max}$：
\begin{equation}
\widetilde w_i
=\frac{\exp(\ell^{(i)}-\ell_{\max})}
{\sum_{j=1}^N\exp(\ell^{(j)}-\ell_{\max})}.
\end{equation}
\begin{block}{归一化权重的含义}
\small
$\ell^{(i)}$ 衡量路径像 Student 的程度；$\ell_{\max}$ 只为数值稳定；$\widetilde w_i$ 是和为 1 的路径比例。
\end{block}
\begin{alertblock}{ESS：有效样本数}
\begin{equation}
\ESS=\left(\sum_i\widetilde w_i^2\right)^{-1}.
\end{equation}
ESS 衡量权重集中程度：全部权重相等时为 $N$；一条路径几乎占据全部权重时接近 $1$。它是常用的近似诊断量，不等于独立样本的精确数量。
\end{alertblock}
\end{frame}

\begin{frame}{步骤3：得到 Student score 并替代 Fake}
查询通道满足
\begin{equation}
X_t=\alpha_tX_{\rm src}+\sigma_t\varepsilon,
\qquad \varepsilon\sim\N(0,I).
\end{equation}
在高斯查询模型下
\begin{equation}
s_\theta(x_t,t)=
\frac{\alpha_t m_{\theta,t}(x_t)-x_t}{\sigma_t^2}.
\end{equation}
所以 YYPF 估计器直接给出
\begin{equation}
\boxed{\widehat s_\theta(x_t,t)=
\frac{\alpha_t\widehat m_{\theta,t}(x_t)-x_t}{\sigma_t^2}}.
\end{equation}
\begin{block}{为什么不再训练 Fake}
Fake 原本就是为了近似 $m_{\theta,t}$。现在均值由当前 Student 路径律在线估计，Tweedie 公式只负责把均值转换成 RTDMD 需要的 score。
\end{block}
\end{frame}

\begin{frame}{步骤3：得到 Student score 并替代 Fake}
\begin{enumerate}\itemsep3pt
\item 当前 Student 参数 $\theta_n$ 生成少步状态 $X_{\theta_n}$，构造 $x_t$。
\item 用随机化后的 \ensuremath{P_{\rm ref}} 生成 $N$ 条 proposal，累计离散 kernel ratio。
\item 归一化权重，计算 $\widehat m_{\theta_n,t}$ 和 $\widehat s_{\theta_n}$。
\item 与 \ensuremath{P_{\rm ref}} score 比较：
\begin{equation}
\Delta s_n=\operatorname{sg}\left[\widehat s_{\theta_n}(x_t,t)-s_\mathrm{ref}(x_t,t)\right].
\end{equation}
\item 用原 RTDMD generator gradient 更新 $\theta_n\to\theta_{n+1}$。
\end{enumerate}
\begin{block}{梯度边界}
建议对 proposal 路径、路径权重、加权均值和 target score 做 stop-gradient。参数更新后重新评价 transition，旧权重不能直接复用。
\end{block}
{\small YYPF 只替换 Fake score 的产生方式，后面的 RTDMD score gap 和更新保持不变。}
\end{frame}

\begin{frame}{数值诊断：权重是否可信}
\begin{block}{有效样本数}
\begin{equation}
\ESS=\left(\sum_{i=1}^N\widetilde w_i^2\right)^{-1}.
\end{equation}
同时记录每一步 log-weight、漂移差范数和权重熵。
\end{block}
\begin{columns}[T,onlytextwidth]
\column{.48\textwidth}
\begin{exampleblock}{ESS 低通常意味着}
\ensuremath{P_{\rm ref}} 与 Student 差异过大、反向窗口太长、粒子数不足，或 $G_k$ 太小导致指数权重爆炸。
\end{exampleblock}
\column{.48\textwidth}
\begin{exampleblock}{可操作的改进}
增加粒子数；分段 rollout 与正确重采样；缩短反向窗口；让 \ensuremath{P_{\rm ref}} 更接近 Student；设计 bridge proposal。
\end{exampleblock}
\end{columns}
\begin{alertblock}{无法直接使用的情况}
支持不重叠、扩散不一致且没有可计算离散核比，或 Student transition 完全不可评价时，A 不能可靠执行。
\end{alertblock}
\end{frame}

\begin{frame}{方案 A 的收益与代价}
\begin{columns}[T,onlytextwidth]
\column{.48\textwidth}
\begin{exampleblock}{收益}
\begin{itemize}\itemsep2pt
\item 不再维护一个滞后的独立 Fake 网络。
\item target score 每轮都由当前 Student 参数定义。
\item \ensuremath{P_{\rm ref}} rollout 可复用成熟随机 diffusion 的采样器。
\item 离散核比给出可执行、可诊断的基线。
\end{itemize}
\end{exampleblock}
\column{.48\textwidth}
\begin{alertblock}{代价}
\begin{itemize}\itemsep2pt
\item 每个查询点要生成多条 \ensuremath{P_{\rm ref}} 路径。
\item 权重可能退化，需监控 ESS。
\item 连续 Girsanov 对共同扩散和绝对连续有严格要求。
\item Student kernel 不可评价时，就无法计算该路径权重。
\end{itemize}
\end{alertblock}
\end{columns}
{\small 方案 A 省去 Fake 网络训练，但增加了路径采样的计算量，并需要控制重要性权重的方差。}
\end{frame}

\begin{frame}{方案 A 的计算流程}
\begin{enumerate}\itemsep6pt
\item RTDMD 需要的是 Student 条件均值 $\E_\theta[X_{\rm src}\mid X_t=x_t]$，Fake 只是这个量的一个独立估计器。
\item 方案 A 用随机 \ensuremath{P_{\rm ref}} 产生候选，再用 Student/\ensuremath{P_{\rm ref}} 路径律比把均值纠正回 Student；离散核比是实现起点，Girsanov 是连续极限。
\item 接入位置就在 Fake score 输出处：$\widehat m\to\widehat s$ 后，后面的 RTDMD score gap 和参数更新保持原样。
\end{enumerate}
\vspace{3mm}
\begin{center}
\begin{tikzpicture}[node distance=4mm, every node/.style={font=\small},
 box/.style={draw=deepblue,fill=softblue,rounded corners,align=center,inner sep=4pt,text width=27mm}]
\node[box] (p) {随机 proposal};
\node[box,right=of p] (w) {路径律校正};
\node[box,right=of w] (m) {条件均值};
\node[box,right=of m] (r) {RTDMD 更新};
\draw[-Latex,thick] (p)--(w);\draw[-Latex,thick] (w)--(m);\draw[-Latex,thick] (m)--(r);
\end{tikzpicture}
\end{center}
\sourceline{来源：用户提供的方案 A 论文、原 Beamer 和 YYPF 模板。}
\end{frame}

\end{document}
\end{document}
\end{document}
